\chapter{Modelos de linguagem como camada de raciocínio}
\label{cap:llm}

\begin{objetivos}
Este capítulo não ensina a construir um modelo de linguagem nem desenvolve a
arquitetura \ing{transformer} em detalhe; ele apresenta o mínimo necessário para
julgar o que essa tecnologia pode e não pode fazer no problema do livro. Ao final,
você será capaz de:
\begin{itemize}[leftmargin=*,itemsep=0.2em]
  \item enunciar a equação da atenção e explicar as duas propriedades que ela trouxe ---
        paralelização e caminho curto entre posições;
  \item explicar o que é uma mistura de especialistas e por que ela reduz o custo de
        inferência;
  \item reconhecer os elementos de um \ing{prompt} com saída estruturada e os critérios
        para validá-la;
  \item reconhecer os modos de falha característicos de modelos de linguagem ---
        alucinação, sensibilidade ao prompt, ausência de calibração --- e as
        salvaguardas correspondentes;
  \item posicionar corretamente o papel de um modelo de linguagem em um sistema de
        monitoramento industrial.
\end{itemize}
\end{objetivos}

\section{A pergunta que sobrou}

Ao final da \cref{parte:evolucao}, o sistema delimita o evento e, no experimento com a
anomalia de ICV, sinaliza corretamente $98{,}2\,\%$ de suas janelas como desconhecidas.
O que ele entrega ao operador é, literalmente, isto:

\begin{center}
\ttfamily
\{ segmento: [14230, 15870], classe: desconhecida, escore: 0.94 \}
\end{center}

O operador precisa então abrir as séries temporais, examinar cada sensor, comparar com o
regime anterior, lembrar de eventos semelhantes e formar uma hipótese. O sistema fez o
trabalho de percepção; todo o trabalho de \emph{interpretação} continua com o humano.

Esta parte do livro investiga se um modelo de linguagem pode assumir parte dessa carga.

\begin{destaque}[O que se está e o que não se está propondo]
\textbf{Não} se propõe substituir os modelos numéricos por um modelo de linguagem. Um
\ing{transformer} de propósito geral é ruim em séries temporais multivariadas
amostradas a cada segundo, e não há razão para esperar o contrário.

Propõe-se acrescentar uma camada \emph{acima} deles. Os modelos numéricos detectam,
segmentam e quantificam; o modelo de linguagem recebe as quantidades resultantes ---
não os dados brutos --- e produz a interpretação. É uma divisão de trabalho por
competência, e o restante desta parte a testa empiricamente.
\end{destaque}

\section{O \ing{transformer}}

\subsection{Atenção}

O mecanismo central, introduzido por \citet{Vaswani2017}, é a atenção por produto escalar
escalado (\cref{eq:atencao}):
\[
  \text{Atenção}(\mat{Q}, \mat{K}, \mat{V}) =
  \operatorname{softmax}\!\left( \frac{\mat{Q}\mat{K}^\top}{\sqrt{d_k}} \right) \mat{V}.
\]

A leitura intuitiva: cada posição da sequência emite uma \termo{consulta}{query}, compara-a
com as \termo{chaves}{keys} de todas as demais posições, e agrega os \termo{valores}{values}
ponderados pela similaridade. O resultado é que qualquer posição pode acessar qualquer
outra em uma única operação --- em contraste com uma LSTM, onde a informação precisa
atravessar todos os passos intermediários.

O fator $\sqrt{d_k}$ evita que o produto escalar cresça com a dimensão e sature a
\ing{softmax}.

\begin{destaque}[Por que isso mudou tudo]
Duas propriedades, e nenhuma delas é ``o modelo entende linguagem''.

\textbf{Paralelização.} Uma LSTM processa a sequência passo a passo, e o passo $t$ depende
do $t-1$. A atenção computa todas as posições simultaneamente. Isso é o que tornou
possível treinar sobre corpora de trilhões de tokens.

\textbf{Caminho curto.} Em uma LSTM, a distância informacional entre duas posições
separadas por $n$ passos é $O(n)$. Na atenção, é $O(1)$. Dependências de longo alcance
deixam de ser um problema arquitetural.
\end{destaque}

\subsection{Da arquitetura ao modelo de linguagem}

Empilhando blocos de atenção com camadas densas, normalização e conexões residuais, e
treinando o resultado para prever o próximo \termo{token}{token} em um corpus muito
grande, obtém-se um modelo de linguagem.

O que emerge desse treinamento --- e permanece objeto de debate legítimo --- é a
capacidade de executar tarefas não vistas explicitamente, quando descritas em linguagem
natural. É essa capacidade que este livro explora, com o ceticismo apropriado.

\section{Mistura de especialistas}

\begin{definicao}{Mistura de especialistas}
Arquitetura em que a camada densa de cada bloco é substituída por um conjunto de
sub-redes especialistas e uma rede de roteamento que seleciona, para cada
\ing{token}, apenas um subconjunto pequeno delas. O modelo tem muitos parâmetros no
total, mas ativa poucos por inferência.
\end{definicao}

O modelo usado nesta parte do livro tem aproximadamente \textbf{397 bilhões de parâmetros
totais} e ativa cerca de \textbf{17 bilhões} por \ing{token} --- uma fração de
aproximadamente $4\,\%$ \citep{Qwen35_2026}.

\begin{destaque}[Por que isso importa para uma aplicação industrial]
Um modelo denso de 397 bilhões de parâmetros exigiria muito mais memória e computação;
sua viabilidade dependeria de hardware, quantização, paralelismo e meta de latência. A
mistura de especialistas reduz o custo por \ing{token} ao ativar apenas parte dos
parâmetros.

Para o problema deste livro, a consequência prática é que a camada de linguagem pode ser
consultada em segundos, contra os \SI{405}{\milli\second} do ensemble numérico e os
segundos ou minutos de histórico das janelas numéricas. Na configuração observada, a
camada de raciocínio não dominou a latência; essa conclusão deve ser reavaliada sob carga
e com a latência de rede da implantação.
\end{destaque}

\section{Como se instrui um modelo de linguagem}

\subsection{Prompt}

\begin{definicao}{Prompt}
Texto de entrada fornecido ao modelo, contendo instruções, contexto, dados e formato de
saída desejado. Diferentemente de um programa, não é executado: é interpretado
estatisticamente.
\end{definicao}

A construção do prompt é a principal alavanca de engenharia disponível, e as decisões
relevantes para esta aplicação são quatro.

\begin{description}[leftmargin=0pt,style=nextline,itemsep=0.4em]

\item[Papel e escopo.] Declarar explicitamente o papel do modelo --- ``você é um
especialista em análise de dados de poços submarinos'' --- e, tão importante quanto,
seus limites --- ``julgue apenas com base nas estatísticas fornecidas''.

\item[Conhecimento de domínio explícito.] Não confiar no que o modelo possa ter
absorvido durante o pré-treinamento. Os perfis de classe do \cref{cap:agente} são
fornecidos no prompt, de forma explícita e verificável.

\item[Formato de saída estruturado.] Exigir JSON com campos definidos. Isso permite
validação automática e integração com o restante do sistema, e é o que separa um
experimento de uma peça de software.

\item[Temperatura baixa.] Reduzir a aleatoriedade da amostragem para privilegiar
consistência sobre criatividade.

\end{description}

\subsection{Saída estruturada}

O agente do \cref{cap:agente} produz, para cada segmento:

\begin{itemize}[leftmargin=*,itemsep=0.2em]
  \item três classes candidatas, ordenadas;
  \item uma justificativa por candidata, referenciando as estatísticas;
  \item um nível de confiança (\ing{High}, \ing{Medium}, \ing{Low});
  \item um nome e uma descrição, caso julgue tratar-se de fenômeno não catalogado.
\end{itemize}

\begin{destaque}[Por que três candidatas e não uma]
Porque uma única resposta força o modelo a comprometer-se, e um modelo que se compromete
sob incerteza produz afirmações confiantes e falsas --- exatamente a patologia que a
\cref{fig:closed-multiclasse} exibiu para o classificador multiclasse.

Três candidatas ordenadas comunicam a estrutura da incerteza. E, como o
\cref{cap:resultados-agente} mostrará, a diferença entre acerto na primeira escolha
($35{,}1\,\%$) e nas três primeiras ($63{,}9\,\%$) é grande o bastante para que essa
decisão de projeto seja o que separa um sistema inútil de um sistema útil.
\end{destaque}

\subsection{Modo de raciocínio}

Modelos recentes podem ser instruídos a gerar uma cadeia de raciocínio intermediária
antes da resposta final. O ganho, em tarefas que exigem comparação de múltiplas
evidências, é consistente; o custo é latência e \ing{tokens}.

Neste trabalho o modo está ativado, com temperatura $0{,}2$, \ing{top-p} $0{,}7$ e limite
de \num{4096} \ing{tokens}.

\begin{destaque}[Cadeia de raciocínio não é explicação]
Distinção importante e frequentemente ignorada. O texto que o modelo gera antes da
resposta é uma sequência que aumenta a probabilidade de acerto; não há garantia de que
ele descreva o processo computacional efetivamente realizado.

Neste sistema, a explicabilidade genuína vem de outro lugar: as estatísticas do
\cref{cap:agente} são calculadas deterministicamente e podem ser verificadas
independentemente. A justificativa do modelo é auditável \emph{contra elas}. É essa
ancoragem --- e não a cadeia de raciocínio --- que torna a saída confiável.
\end{destaque}

\section{Modos de falha}

Um capítulo sobre modelos de linguagem em contexto industrial que não enumere seus modos
de falha seria irresponsável.

\begin{description}[leftmargin=0pt,style=nextline,itemsep=0.5em]

\item[Alucinação.] O modelo produz afirmações plausíveis e falsas com a mesma fluência
com que produz afirmações corretas. Em um contexto onde a saída pode motivar uma
intervenção de milhões de dólares, isso é inaceitável sem salvaguarda.

\emph{Mitigação adotada}: fornecer todas as estatísticas no prompt, restringir a
julgamento sobre elas, exigir justificativa que as referencie, e manter o humano como
decisor.

\item[Sensibilidade ao prompt.] Reformulações semanticamente equivalentes podem alterar
o resultado. Isso compromete a reprodutibilidade.

\emph{Mitigação adotada}: prompt fixo, versionado e reportado integralmente.

\item[Ausência de calibração.] A confiança declarada pelo modelo não corresponde
necessariamente a uma probabilidade. ``High'' não significa $90\,\%$.

\emph{Mitigação adotada}: tratar o nível de confiança como ordinal e reportar a acurácia
empírica separadamente.

\item[Contaminação de dados.] O 3W é público desde 2019, e um modelo treinado em corpus
amplo pode tê-lo visto. Um acerto pode refletir memorização, não raciocínio.

\emph{Mitigação adotada}: nenhuma completamente satisfatória. É uma limitação declarada,
e o \cref{cap:resultados-agente} argumenta a partir dos \emph{padrões} de erro, que
memorização não explicaria bem.

\item[Custo e dependência externa.] Inferência em serviço de terceiros implica custo por
chamada, latência de rede e dependência operacional.

\emph{Mitigação adotada}: acionar o modelo apenas sobre segmentos já detectados como
anômalos --- um punhado por dia, não a cada segundo.

\end{description}

\begin{destaque}[A regra que organiza esta parte do livro]
Um modelo de linguagem \textbf{nunca decide sozinho} em um sistema de monitoramento
industrial. Ele propõe, justifica e ordena hipóteses. A decisão é do operador, e o
sistema deve ser projetado para que essa fronteira seja explícita na interface.
\end{destaque}

\section{Infraestrutura}

O modelo é acessado por um serviço gerenciado de inferência, com os parâmetros
resumidos na \cref{tab:llm-config}.

\begin{table}[htbp]
\centering
\caption{Configuração de inferência do agente.}
\label{tab:llm-config}
\small
\begin{tabular}{ll}
\toprule
\textbf{Parâmetro} & \textbf{Valor} \\
\midrule
Modelo                 & Qwen3.5-397B-A17B (mistura de especialistas) \\
Parâmetros totais      & $\approx 397$ bilhões \\
Parâmetros ativos      & $\approx 17$ bilhões por \ing{token} \\
Serviço                & NVIDIA NIM \\
Modo de raciocínio     & ativado \\
Temperatura            & $0{,}2$ \\
\ing{Top-p}            & $0{,}7$ \\
Limite de \ing{tokens} & \num{4096} \\
Formato de saída       & JSON estruturado \\
\bottomrule
\end{tabular}
\end{table}

\begin{destaque}[Sobre reprodutibilidade]
Modelos servidos por API são atualizados sem aviso, e o mesmo prompt pode produzir saídas
diferentes meses depois. Trabalhos que dependem deles devem reportar data de execução,
identificador exato do modelo e todos os parâmetros de amostragem --- e ainda assim a
reprodutibilidade é inferior à de um modelo local. É um custo real da abordagem, e o
\cref{cap:futuro} discutirá alternativas.
\end{destaque}

\resumocapitulo{%
Modelos de linguagem baseados em \ing{transformer} devem sua eficácia a duas propriedades
do mecanismo de atenção --- paralelização completa e caminho informacional de comprimento
constante entre quaisquer posições --- e não a qualquer compreensão no sentido humano. A
arquitetura de mistura de especialistas permite modelos de centenas de bilhões de
parâmetros com custo de inferência de modelos muito menores: o modelo aqui utilizado tem
$397$ bilhões de parâmetros e ativa $17$ bilhões por \ing{token}. A instrução se dá por
prompt, cujas decisões de projeto relevantes são declarar papel e escopo, fornecer
conhecimento de domínio explicitamente, exigir saída estruturada em JSON e usar
temperatura baixa. O agente produz três candidatas ordenadas, com justificativa e
confiança por candidata --- decisão que se mostrará determinante nos resultados. Os modos
de falha característicos são alucinação, sensibilidade ao prompt, ausência de calibração,
contaminação de dados de treinamento e dependência de serviço externo; cada um exige
salvaguarda explícita. A regra que organiza esta parte: o modelo de linguagem propõe e
justifica, nunca decide.}

\begin{exercicios}
\item Calcule a complexidade computacional da atenção em função do comprimento da
      sequência $n$ e da dimensão $d$. Compare com a de uma LSTM. Para que valores de $n$
      a atenção deixa de ser vantajosa?

\item Um modelo com $397$ bilhões de parâmetros ativa $17$ bilhões por \ing{token}.
      Estime a memória necessária para armazenar os pesos em precisão de 16 bits e a
      necessária para uma inferência. Discuta as implicações para implantação local em
      uma unidade de produção.

\item \textbf{Atividade com serviço externo.} Escreva um prompt que instrua um modelo a
      classificar um segmento anômalo de poço a partir de estatísticas fornecidas, com
      saída em JSON validável. Registre modelo, revisão, data e parâmetros; aplique-o a
      três segmentos e avalie a estabilidade das respostas ao reformular o prompt sem
      alterar seu conteúdo semântico.

\item Discuta a afirmação: ``a cadeia de raciocínio gerada pelo modelo é uma explicação
      de sua decisão''. Apresente um argumento contra e proponha um experimento que
      distinga entre cadeia de raciocínio fiel e cadeia de raciocínio confabulada.

\item Projete um procedimento automático de validação da saída do agente contra as
      estatísticas fornecidas no prompt --- por exemplo, verificar se cada afirmação
      numérica na justificativa corresponde a um valor efetivamente presente na entrada.
      Que fração das alucinações esse procedimento capturaria?

\item O 3W é público desde 2019 e pode ter sido visto durante o pré-treinamento. Proponha
      três estratégias para estimar o grau de contaminação e discuta os limites de cada
      uma.

\item Compare o custo total de propriedade de três arquiteturas: modelo servido por API,
      modelo aberto executado localmente, e modelo aberto ajustado ao domínio. Considere
      custo por inferência, latência, reprodutibilidade, privacidade dos dados e esforço
      de manutenção.
\end{exercicios}

\begin{leituras}
\item \citet{Vaswani2017} --- o artigo que introduz o \ing{transformer}.
\item \citet{Qwen2025} --- os fundamentos da família Qwen3; \citet{Qwen35_2026} --- a
      ficha oficial da versão Qwen3.5-397B-A17B utilizada nos experimentos.
\item \citet{Lopes2026LLM} --- o trabalho que originou esta parte do livro.
\item \citet{Liu2024oilgasLLM} e \citet{Chebbi2025energygpt} --- aplicações de modelos de
      linguagem no setor de petróleo e gás.
\item \citet{Sharma2025xaiagents} --- explicabilidade em sistemas baseados em agentes.
\item \citet{Zhang2025LLMTSFD} --- modelos de linguagem aplicados a detecção de falhas em
      séries temporais.
\end{leituras}
